
\note{There are currently no precompiled versions of the necessary tools
  (see below) available. To use them you must compile them yourself
  from the Rockbox source. See \wikilink{HowToCompile} for more information.
}

In order to install the bootloader, you will need to download the following
files:

\begin{itemize}
  \item An original firmware (OF) from \playerman's website:
    \opt{creativezen}{
      \url{https://support.creative.com/downloads/download.aspx?nDownloadId=10890}
      % FIXME: This link downloads the european variant of the firmware,
      % ZEN_PCFW_L22_1_21_03e.exe. I couldn't find a download link for the
      % US variant ZEN_PCFW_L22_1_21_03.exe
    }

  \item The bootloader:
    \opt{creativezen}{\download{bootloader/creative/zen/bootloader-zen.creative}}

  \item The version of mkzenboot for your operating system from Rockbox's server:
    \download{bootloader/...}\fixme{Not yet available}
\end{itemize}

\subsubsection{Build the Firmware Boot File}

\begin{enumerate}

\item Create a folder called \fname{rbinstall} on your desktop or in your home
  directory, and then extract the OF, the bootloader and mkzenboot into it.

\item Open a new command prompt, and navigate to the ``rbinstall'' folder.

\item Run \fname{mkzenboot}, passing according to the following syntax the name
  of the OF and bootloader files that you downloaded, in addition to a name for
  the output file and the option argument for dualboot.
  You should replace \fname{ZEN\_PCFW\_L22\_1\_21\_03e.exe} by the path to the
  firmware you downloaded from Creative.

  From Windows, enter the following command:
  % The following arguments to the code environment select the correct line out
  % of the code block. This is necessary as using \opt{} in the code block would
  % create unwanted empty lines.
  \begin{code}[firstline=\opt{creativezen}{1},%
              lastline=\opt{creativezen}{1}]
    mkzenboot.exe -i ZEN_PCFW_L22_1_21_03e.exe -b bootloader-zen.creative -o firmware.nk -t dualboot
  \end{code}
  From Linux and Mac OS X enter:%
  \begin{code}[firstline=\opt{creativezen}{1},%
              lastline=\opt{creativezen}{1}]
    ./mkzenboot -i ZEN_PCFW_L22_1_21_03e.exe -b bootloader-zen.creative -o firmware.nk -t dualboot
  \end{code}

\end{enumerate}

\subsubsection{Transfer the Boot File to the \dap}

The last step is to install the bootloader on the \dap{}.
To do this, you need the sendfirm tool in \fname{utils/MTP}. It works both
under Windows and Linux. To build it, change to \fname{utils/MTP} and run
\fname{make} for Linux or \fname{make sendfirm.exe} for Windows. It will produce
a \fname{sendfirm} or \fname{sendfirm.exe} executable. Note that the Windows
version requires \fname{MTP\_DLL.dll} which can be found in \fname{utils/MTP}
(you can also build it from \fname{utils/MTP/MTP\_DLL/}).
\fixme{TODO: provide prebuilt versions and link to it}

\begin{enumerate}
  \item Turn off your \dap{}. You have to use the Reset Button, since Creative OS
    doesn't actually shut down the \dap{} when you ``power off'', but only
    puts it into a low power state and wakes up when you power up again.

  \item Boot into recovery mode by holding the \ButtonPlay{} button on boot.
    It will show a menu in which you should select ``Reload firmware''
    and confirm. After having erased the firmware, you can proceed below.

  \item Run ``sendfirm''. Replace \fname{firmware.nk} by the path to the
    bootloader built in the previous section.
    From Windows:%
    % See above for the code environment arguments explanation.
    \begin{code}[firstline=\opt{creativezen}{1},%
                 lastline=\opt{creativezen}{1}]
      sendfirm.exe firmware.nk
    \end{code}
  From Linux and Max OS X:\
    \begin{code}[firstline=\opt{creativezen}{1},%
                 lastline=\opt{creativezen}{1}]
      ./sendfirm firmware.nk
    \end{code}

\end{enumerate}

Note that after sending the whole file (i.e. progress show 100\%), the tool
will probably display the following error after a timeout:%
\begin{code}[firstline=\opt{creativezen}{1},%
             lastline=\opt{creativezen}{1}]
  Error 2: PTP Layer error 2002: get\_suggested\_storage\_id(): could not get storage id from parent id.
\end{code}
This is expected and you can safely ignore it. Your device should display
``Upgrading firmware'' and perform the upgrade.

Please now go to \reference{ref:finish_install} to complete the installation procedure.
